\section{Introduction}
\label{sec:intro}

\begin{figure}[hbt!]
    \includegraphics[width=\linewidth]{sections/rubric_main.pdf}
    \vspace{-1cm}
    \caption{This diagram provides a walkthrough of the experimental setup as shown for automatic essay scoring.
    It represents comparisons made between human-autorater agreements $\tau$ across holistic rubrics (left), in which all criteria are applied together in a single overall judgment, or analytic rubrics (right), in which criteria are evaluated separately, resulting in multiple scores. The original rubrics are given to a human and an autorater, while edited rubrics are given to autoraters only. Arrows in bold between $\tau_1$, $\tau_2$ and $\tau_3$, $\tau_4$ represent comparisons for which statistical significance can be calculated. $\Delta\textnormal{rater}$ represents comparisons where the type of rater is changed while the type of rubric remains constant, while $\Delta\textnormal{rubric}$ represents comparisons where the type of rubric is changed and the type of rater remains constant.}
    \label{fig:experiment_design}
\end{figure}

Autoraters, or LLM-as-judges, have been used as an alternative to human annotation due to their scalability, cost, and time effectiveness. Their effectiveness is typically validated through agreement with human annotation. 
Borrowing from education literature, a \textit{rubric} is defined as having ``\textit{coherent sets of criteria}" and ``\textit{descriptions of levels of performance for these criteria}" \cite{brookhart2013create}. This describes the scoring guidelines and instructions provided to any rater, whether human or automated, which is also referred to as part of a prompt as described in autorater literature.
Ideally, both human raters and autoraters would receive equivalent evaluation rubrics that accurately measure the same construct with reliable certainty.
% Ideally, both human annotators and autorators should receive equivalent evaluation rubrics that measure the same construct and elicit the equally reliable judgments. 
However, equivalence does not necessitate identical presentation. 
\citet{Wu_Quinn_2017} show that expert and non-expert human raters may require different levels of instruction specificity, for instance, specifying tools and providing concrete examples improves accuracy specifically when raters lack task-relevant knowledge. 

In addition to identical presentation across different types of raters, humans and autoraters are sensitive to variations in instruction presentations within the same type of rater. For humans, this sensitivity includes the interpretation of the task during crowdsourcing \citep{10.1145/2818048.2820016} and instruction specification, where increasing instruction specification increases accuracy on the task \citep{Wu_Quinn_2017}. For autoraters, position bias (where the position of the evaluated text within the prompt will influence the autorater evaluation) and verbosity bias (where autoraters prefer more verbose texts) are only a few of the known sensitivities \citep{10.5555/3666122.3668142}. Autoraters are also sensitive to rubric variations such as formatting choices \cite{sclar2024quantifying} and example ordering \cite{lu-etal-2022-fantastically}. Another example of instruction presentation involves decomposing criteria into multiple sub-criteria. Previous work has successfully used decomposition to improve LLM performance, whether by having autoraters decompose evaluation criteria into sub-tasks \citep{saha-etal-2024-branch} or by having humans decompose complex questions into simpler sub-questions for models
\citep{patel-etal-2022-question}.

Understanding whether rubric modifications produce statistically significant shifts in agreement is essential for practitioners who aim to deploy autoraters as evaluation tools. This work studies human-autorater agreement in two domains: automatic essay scoring (AES) and instruction-following (IF). We examine statistically how rubric presentation and broader rubric modifications affect human-autorater agreement on subjective evaluation tasks, as well as empirically examining whether decomposing general holistic judgments that ask for a single high-level criteria into more granular sub-criteria, originally designed for human evaluators, can be an effective approach for improving human-judge agreement. The findings indicate that instructions optimized for autoraters tend to improve agreement with human ratings when autoraters receive machine-optimized instructions and humans receive the original set of instructions.
Conversely, giving autoraters simpler prompts does not guarantee higher agreement with human ratings.
These results indicate that 1) rubric edits providing representative examples along with contextual information increased human-autorater agreement as well as autorater self-agreement, 2) higher criterion complexity and conservative aggregation methods tended to decrease human-autorater agreement, 3) reducing confirmation bias tends to significantly increase human-autorater agreement and 4) high human inter-rater agreement leads to significantly higher human-autorater agreement.
These findings from the automatic essay scoring and instruction-following evaluation domains suggest that practitioners should carefully analyze domain-specific performance and modify rubrics to move towards high human-autorater agreement. 
% \alfredo{The results indicate that rubric edits providing representative examples and additional context increased human-autorater agreement, while higher rubric complexity and conservative aggregation methods tended to decrease it. The findings from the automatic essay scoring and instruction-following evaluation domains suggest that practitioners should carefully analyze domain-specific performance and modify rubrics to move towards high human-autorater agreement.} \jessica{wait can we rewrite this somehow to not be exactly the same as the abstract}